\begin{table}[H]
\centering
\caption{Endline Valence Contrasts Relative to Apolitical China Content}
\label{tab:valence_vs_apolitical_t1}
\footnotesize
\begin{threeparttable}
\begin{tabular}{llcccc}
\toprule
Outcome & Contrast & Estimate & Std. Error & p-value & N \\
\midrule
Censorship support & Pooled Anti-China vs Apolitical & -0.107*** & 0.027 & <0.001 & 4,348 \\
Censorship support & Pooled Pro-China vs Apolitical & 0.008 & 0.027 & 0.755 & 4,348 \\
China-West thermometer gap & Pooled Anti-China vs Apolitical & -2.076*** & 0.487 & <0.001 & 4,348 \\
China-West thermometer gap & Pooled Pro-China vs Apolitical & -1.904*** & 0.481 & <0.001 & 4,348 \\
Nationalism & Pooled Anti-China vs Apolitical & 0.119*** & 0.026 & <0.001 & 4,348 \\
Nationalism & Pooled Pro-China vs Apolitical & 0.052* & 0.026 & 0.046 & 4,348 \\
Regime support & Pooled Anti-China vs Apolitical & -0.284*** & 0.023 & <0.001 & 4,348 \\
Regime support & Pooled Pro-China vs Apolitical & 0.189*** & 0.023 & <0.001 & 4,348 \\
Trust in foreign media & Pooled Anti-China vs Apolitical & -0.045 & 0.029 & 0.116 & 4,348 \\
Trust in foreign media & Pooled Pro-China vs Apolitical & 0.140*** & 0.029 & <0.001 & 4,348 \\
\bottomrule
\end{tabular}
\begin{tablenotes}[flushleft]
\footnotesize
\item Note: Entries report pooled Pro-China and pooled Anti-China contrasts relative to the apolitical-China group, estimated on the five non-control groups only. All specifications include the baseline value of the dependent variable and randomization-block fixed effects. This table is designed to map as closely as possible to the registered valence comparisons that used apolitical content as the neutral benchmark.
\end{tablenotes}
\end{threeparttable}
\end{table}
